\documentclass[11pt]{article}

\usepackage[utf8]{inputenc}
\usepackage[T1]{fontenc}
\usepackage[spanish]{babel}
\usepackage{amsmath,amssymb,amsthm,mathtools}
\usepackage{booktabs}
\usepackage{array}
\usepackage{longtable}
\usepackage{geometry}
\usepackage{hyperref}
\usepackage{enumitem}
\usepackage{graphicx}
\usepackage{xcolor}

\geometry{margin=2.2cm}
\setlength{\parskip}{0.45em}
\setlength{\parindent}{0pt}
\hypersetup{colorlinks=true,linkcolor=blue,urlcolor=blue}

\title{Gu\'ia de Estudio Completa de Control\\
\large Basada en las clases de estabilidad, respuesta temporal, PID y sintonizaci\'on}
\author{Sergio E. Ropero}
\date{\today}

\begin{document}

\maketitle
\tableofcontents
\newpage

\section{C\'omo usar este documento}
Este documento est\'a pensado como una gu\'ia de estudio completa para repasar los temas vistos en clase. La idea no es solo listar f\'ormulas, sino explicar de d\'onde salen, qu\'e significan y c\'omo se usan en ejercicios.

Los bloques principales del documento son:
\begin{enumerate}[label=\arabic*.]
\item Modelado, linealizaci\'on y funci\'on de transferencia.
\item Polos, ceros, estabilidad y criterio de Routh-Hurwitz.
\item Respuesta temporal y m\'etricas de sistemas de primer y segundo orden.
\item Terminolog\'ia de lazos de control y respuesta a escal\'on de procesos.
\item Controladores PID: teor\'ia, estructuras y efecto de cada t\'ermino.
\item M\'etodos de sintonizaci\'on: Ziegler-Nichols, Cohen-Coon y s\'intesis directa.
\item Ejemplos resueltos paso a paso.
\end{enumerate}

En varios ejemplos las diapositivas muestran solo el enunciado o parte del desarrollo. En esos casos, el procedimiento que aparece aqu\'i fue reconstruido siguiendo la misma metodolog\'ia explicada en clase.

\section{Panorama general del curso}
La l\'ogica general del curso de control suele seguir esta secuencia:
\begin{enumerate}[label=\arabic*.]
\item Se obtiene un modelo matem\'atico del sistema din\'amico.
\item Si el modelo es no lineal, se linealiza alrededor de un punto de equilibrio.
\item Se escribe el sistema en espacio de estados o como funci\'on de transferencia.
\item Se analizan polos, ceros, estabilidad y respuesta temporal.
\item Se dise\~na un controlador para lograr un comportamiento deseado.
\item Se valida el dise\~no en simulaci\'on y luego, si aplica, en implementaci\'on real.
\end{enumerate}

Ese flujo es importante porque casi todos los ejercicios de la materia son una variaci\'on de ese mismo esquema.

\section{Modelado y funci\'on de transferencia}
\subsection{Modelo no lineal y linealizaci\'on}
En forma general, un sistema no lineal puede escribirse como
\begin{align}
\dot{x} &= f(x,u),\\
y &= h(x,u),
\end{align}
donde:
\begin{itemize}
\item $x \in \mathbb{R}^n$ es el vector de estados,
\item $u \in \mathbb{R}^m$ es la entrada,
\item $y \in \mathbb{R}^p$ es la salida.
\end{itemize}

Si se tiene un punto de equilibrio $(x^\star,u^\star)$, se define
\[
\hat{x}=x-x^\star,\qquad \hat{u}=u-u^\star,\qquad \hat{y}=y-y^\star
\]
y el sistema linealizado queda
\begin{align}
\dot{\hat{x}} &= A\hat{x}+B\hat{u},\\
\hat{y} &= C\hat{x}+D\hat{u},
\end{align}
con
\[
A=\left.\frac{\partial f}{\partial x}\right|_\star,\quad
B=\left.\frac{\partial f}{\partial u}\right|_\star,\quad
C=\left.\frac{\partial h}{\partial x}\right|_\star,\quad
D=\left.\frac{\partial h}{\partial u}\right|_\star.
\]

\subsection{Funci\'on de transferencia desde espacio de estados}
Para un sistema lineal e invariante en el tiempo, la funci\'on de transferencia se obtiene como
\begin{equation}
G(s)=\frac{Y(s)}{U(s)}=C(sI-A)^{-1}B+D.
\end{equation}

Aqu\'i aparece la matriz
\[
\Phi(s)=(sI-A)^{-1},
\]
y en tiempo la matriz de transici\'on de estado es
\[
\Phi(t)=\mathcal{L}^{-1}\{(sI-A)^{-1}\}=e^{At}.
\]

\subsection{Polos, ceros y ganancia}
Si la funci\'on de transferencia se escribe como
\begin{equation}
G(s)=\frac{N(s)}{D(s)},
\end{equation}
entonces:
\begin{itemize}
\item los \textbf{ceros} son las ra\'ices de $N(s)=0$,
\item los \textbf{polos} son las ra\'ices de $D(s)=0$,
\item la \textbf{ganancia DC} es $G(0)$ si existe.
\end{itemize}

Los polos son especialmente importantes porque determinan estabilidad, rapidez, oscilaci\'on y amortiguamiento.

\subsection{Ejemplo b\'asico: Tower Copter}
Si el sistema est\'a descrito por
\begin{equation}
2.4712\dot{y}(t)+y(t)=0.32029\,u(t),
\end{equation}
entonces aplicando transformada de Laplace con condiciones iniciales nulas:
\[
2.4712\,sY(s)+Y(s)=0.32029\,U(s).
\]
Agrupando $Y(s)$:
\[
(2.4712s+1)Y(s)=0.32029\,U(s),
\]
y por tanto
\begin{equation}
G(s)=\frac{Y(s)}{U(s)}=\frac{0.32029}{2.4712s+1}.
\end{equation}

Este sistema:
\begin{itemize}
\item no tiene ceros,
\item tiene un polo en $s=-1/2.4712=-0.4047$,
\item es estable porque su polo est\'a en el semiplano izquierdo.
\end{itemize}

\section{Respuesta temporal}
\subsection{Sistema de primer orden}
La forma can\'onica de un sistema de primer orden es
\begin{equation}
G(s)=\frac{K}{\tau s+1}.
\end{equation}

Su respuesta al escal\'on unitario es
\begin{equation}
y(t)=K\left(1-e^{-t/\tau}\right),\qquad t\ge 0.
\end{equation}

Interpretaci\'on de los par\'ametros:
\begin{itemize}
\item $K$ es la ganancia est\'atica o ganancia DC.
\item $\tau$ es la constante de tiempo.
\end{itemize}

Reglas r\'apidas:
\begin{itemize}
\item En $t=\tau$, la respuesta alcanza aproximadamente el $63.2\%$ del valor final.
\item El tiempo de establecimiento al $2\%$ se aproxima por
\[
t_{ss}\approx 4\tau.
\]
\item El tiempo de establecimiento al $5\%$ se aproxima por
\[
t_{ss}\approx 3\tau.
\]
\end{itemize}

\subsection{Ejemplo guiado: respuesta de primer orden del Tower Copter}
Tomemos la funci\'on de transferencia
\[
G(s)=\frac{0.32029}{2.4712s+1}.
\]

\textbf{Paso 1: identificar la forma can\'onica.}

La comparamos con
\[
G(s)=\frac{K}{\tau s+1}.
\]
De inmediato se reconoce que
\[
K=0.32029,\qquad \tau=2.4712~\text{s}.
\]

\textbf{Por qu\'e se hace esto.}

Porque al reconocer la forma can\'onica de primer orden ya sabemos c\'omo luce la respuesta temporal y podemos sacar conclusiones sin volver a derivar todo desde cero.

\textbf{Paso 2: calcular la respuesta al escal\'on unitario.}

Para una entrada escal\'on unitario,
\[
R(s)=\frac{1}{s}.
\]
Entonces
\[
Y(s)=G(s)R(s)=\frac{0.32029}{s(2.4712s+1)}.
\]
La transformada inversa conocida de un sistema de primer orden da
\[
y(t)=0.32029\left(1-e^{-t/2.4712}\right),\qquad t\ge 0.
\]

\textbf{Paso 3: interpretar la expresi\'on.}

Esta respuesta muestra que:
\begin{enumerate}[label=\arabic*.]
\item la salida crece de manera mon\'otona;
\item no hay oscilaciones;
\item el valor final es
\[
\lim_{t\to\infty}y(t)=0.32029;
\]
\item el sistema es estable, porque el exponencial decae.
\end{enumerate}

\textbf{Paso 4: usar la constante de tiempo.}

En $t=\tau=2.4712$ s:
\[
y(\tau)=0.32029\left(1-e^{-1}\right)\approx 0.32029(0.632)\approx 0.2024.
\]
Eso confirma la regla del $63.2\%$: en una constante de tiempo ya recorri\'o el $63.2\%$ del cambio total.

\textbf{Paso 5: estimar el tiempo de establecimiento.}

Usando la regla del $2\%$:
\[
t_{ss}\approx 4\tau=4(2.4712)=9.8848~\text{s}.
\]

\textbf{Conclusi\'on del ejemplo.}

Este ejemplo ilustra muy bien por qu\'e en sistemas de primer orden la constante de tiempo es la cantidad m\'as importante: de ella dependen casi todas las conclusiones temporales b\'asicas.

\subsection{Sistema de segundo orden}
La forma can\'onica m\'as importante es
\begin{equation}
G(s)=\frac{\omega_n^2}{s^2+2\zeta \omega_n s+\omega_n^2},
\end{equation}
donde:
\begin{itemize}
\item $\omega_n$ es la frecuencia natural,
\item $\zeta$ es el coeficiente de amortiguamiento.
\end{itemize}

Los polos son
\begin{equation}
s_{1,2}=-\zeta \omega_n \pm j\omega_n\sqrt{1-\zeta^2}.
\end{equation}

\subsection{Clasificaci\'on seg\'un $\zeta$}
\begin{itemize}
\item $\zeta>1$: sobreamortiguado.
\item $\zeta=1$: cr\'iticamente amortiguado.
\item $0<\zeta<1$: subamortiguado.
\item $\zeta=0$: no amortiguado.
\item $\zeta<0$: inestable.
\end{itemize}

\subsection{M\'etricas de desempe\~no para segundo orden subamortiguado}
Definiendo
\[
\omega_d=\omega_n\sqrt{1-\zeta^2},
\]
se tienen las f\'ormulas m\'as usadas:

\begin{align}
t_p &= \frac{\pi}{\omega_d},\\
M_p &= e^{-\pi\zeta/\sqrt{1-\zeta^2}}\times 100\%,\\
t_{ss}(2\%) &\approx \frac{4}{\zeta\omega_n}.
\end{align}

Una expresi\'on com\'un para el tiempo de subida es
\begin{equation}
t_r=\frac{\pi-\phi}{\omega_d},
\qquad
\phi=\arctan\!\left(\frac{\sqrt{1-\zeta^2}}{\zeta}\right)=\arccos(\zeta).
\end{equation}

\subsection{Ejemplo guiado: m\'etricas de un segundo orden}
Consideremos
\[
G(s)=\frac{16}{s^2+4s+16}.
\]

\textbf{Paso 1: llevar el sistema a la forma est\'andar.}

La forma est\'andar es
\[
G(s)=\frac{\omega_n^2}{s^2+2\zeta\omega_n s+\omega_n^2}.
\]
Comparando coeficientes:
\[
\omega_n^2=16 \quad\Rightarrow\quad \omega_n=4,
\]
y
\[
2\zeta\omega_n=4
\quad\Rightarrow\quad
2\zeta(4)=4
\quad\Rightarrow\quad
\zeta=0.5.
\]

\textbf{Por qu\'e se hace esto.}

Porque las f\'ormulas de desempe\~no dependen de $\zeta$ y $\omega_n$, no del polinomio en bruto.

\textbf{Paso 2: calcular la frecuencia amortiguada.}

\[
\omega_d=\omega_n\sqrt{1-\zeta^2}
=4\sqrt{1-0.25}
=4\sqrt{0.75}
\approx 3.4641.
\]

\textbf{Paso 3: calcular el tiempo de pico.}

\[
t_p=\frac{\pi}{\omega_d}
\approx \frac{3.1416}{3.4641}
\approx 0.9069~\text{s}.
\]

\textbf{Paso 4: calcular el sobreimpulso.}

\[
M_p=e^{-\pi\zeta/\sqrt{1-\zeta^2}}100\%
=e^{-\pi(0.5)/\sqrt{0.75}}100\%
\approx 16.30\%.
\]

\textbf{Paso 5: calcular el tiempo de establecimiento.}

\[
t_{ss}\approx \frac{4}{\zeta\omega_n}
=\frac{4}{0.5(4)}
=2~\text{s}.
\]

\textbf{Paso 6: calcular el tiempo de subida.}

Primero:
\[
\phi=\arccos(\zeta)=\arccos(0.5)=\frac{\pi}{3}.
\]
Luego:
\[
t_r=\frac{\pi-\phi}{\omega_d}
=\frac{\pi-\pi/3}{3.4641}
\approx 0.6046~\text{s}.
\]

\textbf{Conclusi\'on del ejemplo.}

Como $0<\zeta<1$, el sistema es subamortiguado. Eso explica por qu\'e presenta sobreimpulso y tiempo pico finito. Este ejemplo es bueno para practicar porque obliga a pasar del polinomio original a los par\'ametros f\'isicos del segundo orden.

\subsection{Qu\'e significa cada m\'etrica?}
\begin{itemize}
\item \textbf{Tiempo de subida} $t_r$: qu\'e tan r\'apido la salida empieza a alcanzar la referencia.
\item \textbf{Tiempo de pico} $t_p$: instante del primer m\'aximo.
\item \textbf{Sobreimpulso} $M_p$: qu\'e tanto se pasa del valor final.
\item \textbf{Tiempo de establecimiento} $t_{ss}$: cu\'ando la respuesta entra y se queda dentro de una banda alrededor del valor final.
\end{itemize}

\section{Estabilidad}
\subsection{Idea general}
En sistemas LTI, la estabilidad est\'a fuertemente asociada a la ubicaci\'on de los polos.

Regla base:
\begin{itemize}
\item Si todos los polos tienen parte real negativa, el sistema es asint\'oticamente estable.
\item Si alg\'un polo tiene parte real positiva, el sistema es inestable.
\item Si hay polos sobre el eje imaginario y no est\'an amortiguados por otros efectos, el caso es marginal o cr\'itico y hay que analizarlo con cuidado.
\end{itemize}

\subsection{Condici\'on necesaria}
Si el polinomio caracter\'istico es
\[
a_ns^n+a_{n-1}s^{n-1}+\cdots+a_1s+a_0=0,
\]
una condici\'on necesaria para estabilidad es que todos los coeficientes sean del mismo signo.

Ojo: eso es \textbf{necesario}, pero no siempre \textbf{suficiente}.

\subsection{Criterio de Routh-Hurwitz}
Este criterio permite determinar cu\'antos polos est\'an en el semiplano derecho sin calcularlos expl\'icitamente.

Para un polinomio
\[
a_ns^n+a_{n-1}s^{n-1}+\cdots+a_1s+a_0=0,
\]
se construye la tabla de Routh. La estabilidad asint\'otica se garantiza si y solo si todos los elementos de la primera columna son positivos.

\subsubsection{Ejemplo de tabla para un c\'ubico}
Para
\[
a_3s^3+a_2s^2+a_1s+a_0=0,
\]
la tabla es
\[
\begin{array}{c|cc}
s^3 & a_3 & a_1\\
s^2 & a_2 & a_0\\
s^1 & \dfrac{a_2a_1-a_3a_0}{a_2} & 0\\
s^0 & a_0 & 0
\end{array}
\]
y las condiciones de estabilidad son:
\[
a_3>0,\quad a_2>0,\quad a_0>0,\quad a_2a_1-a_3a_0>0.
\]

\subsection{Ejemplo resuelto: rango de $K_c$ para estabilidad}
En las diapositivas aparece el proceso
\begin{equation}
G_p(s)=\frac{1}{(s+1)(0.5s+1)(0.2s+1)}
\end{equation}
con realimentaci\'on unitaria y un controlador proporcional $K_c$.

La ecuaci\'on caracter\'istica en lazo cerrado es
\[
1+K_cG_p(s)=0.
\]

Como
\[
(s+1)(0.5s+1)(0.2s+1)=0.1s^3+0.8s^2+1.7s+1,
\]
la ecuaci\'on caracter\'istica queda
\begin{equation}
0.1s^3+0.8s^2+1.7s+(1+K_c)=0.
\label{eq:cubic_example}
\end{equation}

La tabla de Routh es
\[
\begin{array}{c|cc}
s^3 & 0.1 & 1.7\\
s^2 & 0.8 & 1+K_c\\
s^1 & \dfrac{0.8(1.7)-0.1(1+K_c)}{0.8} & 0\\
s^0 & 1+K_c & 0
\end{array}
\]

Las condiciones de estabilidad son:
\begin{align}
1+K_c &>0,\\
\frac{0.8(1.7)-0.1(1+K_c)}{0.8} &>0.
\end{align}

La primera da
\[
K_c>-1.
\]

La segunda da
\[
1.36-0.1-0.1K_c>0
\]
\[
1.26-0.1K_c>0
\]
\[
K_c<12.6.
\]

Por tanto,
\begin{equation}
-1<K_c<12.6.
\end{equation}

Como en la pr\'actica normalmente se trabaja con ganancias positivas,
\begin{equation}
0<K_c<12.6.
\end{equation}

\subsection{Ganancia cr\'itica y per\'iodo cr\'itico}
Para hallar la ganancia cr\'itica $K_u$ y el per\'iodo cr\'itico $P_u$, se sustituye $s=j\omega$ en \eqref{eq:cubic_example}:
\[
0.1(j\omega)^3+0.8(j\omega)^2+1.7(j\omega)+(1+K_c)=0.
\]

Separando parte real e imaginaria:
\begin{align}
-0.8\omega^2+1+K_c &= 0,\\
-0.1\omega^3+1.7\omega &= 0.
\end{align}

De la parte imaginaria:
\[
\omega(-0.1\omega^2+1.7)=0.
\]
Descartando $\omega=0$:
\[
0.1\omega^2=1.7
\quad\Rightarrow\quad
\omega^2=17
\quad\Rightarrow\quad
\omega_u=\sqrt{17}.
\]

Entonces
\begin{equation}
P_u=\frac{2\pi}{\omega_u}=\frac{2\pi}{\sqrt{17}}\approx 1.5239~\text{s}.
\end{equation}

De la parte real:
\[
K_u=0.8\omega_u^2-1=0.8(17)-1=12.6.
\]

Por tanto:
\begin{equation}
K_u=12.6,\qquad P_u\approx 1.5239~\text{s}.
\end{equation}

\section{Sistemas con tiempo muerto y aproximaci\'on de Pad\'e}
Un tiempo muerto o retardo se representa como
\[
e^{-Ls}.
\]

El problema es que no es una funci\'on racional, y muchas t\'ecnicas cl\'asicas trabajan mejor con polinomios. Por eso se usa la aproximaci\'on de Pad\'e. La de primer orden es
\begin{equation}
e^{-Ls}\approx \frac{1-\frac{L}{2}s}{1+\frac{L}{2}s}.
\end{equation}

Esto permite llevar el problema a una forma racional para estudiar estabilidad con Routh-Hurwitz.

\subsection{Ejemplo guiado: estabilidad con tiempo muerto usando Pad\'e}
Supongamos que se tiene un sistema con realimentaci\'on unitaria y planta
\[
G(s)=\frac{K e^{-Ls}}{s+1}.
\]
Queremos hallar un rango de $K$ que garantice estabilidad.

\textbf{Paso 1: aproximar el retardo.}

Usamos Pad\'e de primer orden:
\[
e^{-Ls}\approx \frac{1-\frac{L}{2}s}{1+\frac{L}{2}s}.
\]
Entonces
\[
G(s)\approx \frac{K\left(1-\frac{L}{2}s\right)}{(s+1)\left(1+\frac{L}{2}s\right)}.
\]

\textbf{Por qu\'e se hace esto.}

Porque Routh-Hurwitz necesita una ecuaci\'on caracter\'istica polinomial. El retardo puro no permite eso directamente.

\textbf{Paso 2: escribir la ecuaci\'on caracter\'istica.}

Con realimentaci\'on unitaria:
\[
1+G(s)=0.
\]

Sustituyendo:
\[
1+\frac{K\left(1-\frac{L}{2}s\right)}{(s+1)\left(1+\frac{L}{2}s\right)}=0.
\]

Multiplicando por el denominador:
\[
(s+1)\left(1+\frac{L}{2}s\right)+K\left(1-\frac{L}{2}s\right)=0.
\]

\textbf{Paso 3: expandir.}

\[
\frac{L}{2}s^2+\left(1+\frac{L}{2}\right)s+1+K-\frac{KL}{2}s=0.
\]
Agrupando t\'erminos:
\begin{equation}
\frac{L}{2}s^2+\left(1+\frac{L}{2}-\frac{KL}{2}\right)s+(1+K)=0.
\end{equation}

\textbf{Paso 4: imponer estabilidad.}

Para un segundo orden, basta exigir que todos los coeficientes sean positivos:
\begin{align}
\frac{L}{2} &>0,\\
1+\frac{L}{2}-\frac{KL}{2} &>0,\\
1+K &>0.
\end{align}

La tercera condici\'on da
\[
K>-1.
\]
La segunda da
\[
1+\frac{L}{2}-\frac{KL}{2}>0
\]
\[
2+L-KL>0
\]
\[
K<1+\frac{2}{L}.
\]

\textbf{Resultado final.}

\[
-1<K<1+\frac{2}{L}.
\]

Si en la pr\'actica se exige ganancia positiva:
\[
0<K<1+\frac{2}{L}.
\]

\textbf{Conclusi\'on del ejemplo.}

El retardo hace m\'as dif\'icil la estabilizaci\'on. Mientras mayor es $L$, menor es el valor m\'aximo permisible de la ganancia.

\section{Estabilidad robusta y criterio de Kharitonov}
Cuando los coeficientes del polinomio no son exactos sino que pertenecen a intervalos, ya no basta analizar un \'unico polinomio. Kharitonov da una forma de estudiar estabilidad robusta sin probar infinitos casos.

La idea importante para estudiar es:
\begin{itemize}
\item si los coeficientes del polinomio var\'ian dentro de intervalos,
\item la estabilidad de toda la familia puede inferirse verificando cuatro polinomios especiales de Kharitonov.
\end{itemize}

No siempre piden desarrollar esto en detalle, pero s\'i conviene saber qu\'e resuelve: \textbf{estabilidad frente a incertidumbre param\'etrica}.

\section{Lazos de control y terminolog\'ia de procesos}
\subsection{Estructura general de un lazo cerrado}
En procesos industriales aparece el siguiente esquema conceptual:
\[
R(s)\rightarrow \text{controlador} \rightarrow \text{actuador} \rightarrow \text{proceso} \rightarrow \text{sensor} \rightarrow Y(s).
\]

Con mayor detalle:
\begin{itemize}
\item $R(s)$: referencia o set-point.
\item $E(s)=R(s)-Y(s)$: error.
\item $G_c(s)$: controlador.
\item $G_v(s)$: elemento final de control o actuador.
\item $G_p(s)$: planta o proceso.
\item $G_t(s)$, $G_s(s)$: transmisor y sensor.
\item $D(s)$: perturbaci\'on.
\item $N(s)$: ruido de medici\'on.
\end{itemize}

\subsection{T\'erminos m\'as usados}
\begin{itemize}
\item \textbf{PV} (Process Variable): variable del proceso.
\item \textbf{SP} (Set Point): referencia deseada.
\item \textbf{MV} (Manipulated Variable): variable manipulada.
\item \textbf{CV} (Controlled Variable): variable controlada.
\item \textbf{Offset}: error permanente entre PV y SP.
\end{itemize}

\subsection{Ejemplos t\'ipicos}
\begin{itemize}
\item En control de temperatura: PV = temperatura, MV = flujo de combustible o potencia de calefacci\'on.
\item En control de nivel: PV = nivel del tanque, MV = apertura de v\'alvula o caudal de entrada.
\end{itemize}

\section{An\'alisis de respuesta a escal\'on en procesos}
\subsection{Modelo de primer orden con tiempo muerto}
Muchos procesos se aproximan por
\begin{equation}
G_p(s)=\frac{Ke^{-Ls}}{Ts+1},
\end{equation}
donde:
\begin{itemize}
\item $K$ es la ganancia del proceso,
\item $L$ es el tiempo muerto,
\item $T$ es la constante de tiempo.
\end{itemize}

\subsection{M\'etodo del 63\%}
Este es uno de los m\'etodos m\'as importantes de identificaci\'on a partir de respuesta escal\'on.

Pasos:
\begin{enumerate}[label=\arabic*.]
\item Aplicar un cambio tipo escal\'on a la entrada del proceso.
\item Encontrar el punto de la curva con mayor pendiente.
\item Trazar la tangente en ese punto.
\item Hallar la intersecci\'on de la tangente con el valor inicial de la salida. Ese tiempo define $L$.
\item Hallar cu\'ando la salida llega al $63\%$ del cambio total. La diferencia con $L$ da $T$.
\item Calcular la ganancia del proceso como
\begin{equation}
K=\frac{\Delta y}{\Delta u}.
\end{equation}
\end{enumerate}

\subsection{Caso de procesos integradores}
Para un proceso integrador no tiene mucho sentido hablar de ganancia DC finita. En su lugar aparece la \textbf{velocidad de ganancia}
\begin{equation}
K_v=\frac{\Delta y}{\Delta u\,L}.
\end{equation}

La interpretaci\'on es: qu\'e tan r\'apido crece la salida por unidad de cambio en la entrada.

\subsection{Ejemplo aproximado con el proceso de tercer orden}
Tomemos otra vez
\[
G_p(s)=\frac{1}{(s+1)(0.5s+1)(0.2s+1)}.
\]

Usando num\'ericamente el m\'etodo de la tangente en la respuesta al escal\'on, se obtiene aproximadamente:
\begin{equation}
K\approx 1,\qquad L\approx 0.359~\text{s},\qquad T\approx 2.103~\text{s}.
\end{equation}

Estos valores no son exactos del proceso real, sino una aproximaci\'on FOPDT para usar m\'etodos de sintonizaci\'on basados en respuesta temporal.

\subsection{Ejemplo guiado: c\'omo se calculan $L$ y $T$ con la tangente}
Vamos a explicar con detalle c\'omo se llega a los valores
\[
L\approx 0.359,\qquad T\approx 2.103
\]
para el proceso
\[
G_p(s)=\frac{1}{(s+1)(0.5s+1)(0.2s+1)}.
\]

\textbf{Paso 1: obtener la respuesta al escal\'on del proceso.}

Se simula la salida $y(t)$ ante una entrada escal\'on unitario. Como la ganancia DC del proceso es $1$, el valor final de la respuesta es tambi\'en $1$.

\textbf{Paso 2: encontrar el punto de m\'axima pendiente.}

Num\'ericamente se encuentra un punto de inflexi\'on aproximado en
\[
t_i\approx 0.9435,\qquad y_i\approx 0.2779,
\]
con pendiente
\[
m\approx 0.4755.
\]

\textbf{Por qu\'e se usa este punto.}

Porque la tangente en el punto de inflexi\'on es la que mejor aproxima la parte central de la curva y permite separar el efecto del retardo y de la constante de tiempo.

\textbf{Paso 3: escribir la ecuaci\'on de la tangente.}

La recta tangente en ese punto es
\[
y_{\text{tan}}(t)=m(t-t_i)+y_i.
\]

\textbf{Paso 4: hallar $L$.}

Por construcci\'on, el tiempo muerto aparente $L$ se obtiene haciendo que la tangente corte el valor inicial de la salida. Como la salida arranca en $0$:
\[
0=m(t-t_i)+y_i.
\]
Despejando:
\[
t=t_i-\frac{y_i}{m}.
\]
Entonces
\[
L=t_i-\frac{y_i}{m}
\approx 0.9435-\frac{0.2779}{0.4755}
\approx 0.359~\text{s}.
\]

\textbf{Paso 5: hallar el tiempo donde la tangente corta el valor final.}

Como el valor final es $1$:
\[
1=m(t-t_i)+y_i.
\]
Despejando:
\[
t=t_i+\frac{1-y_i}{m}.
\]
Num\'ericamente:
\[
t_f\approx 0.9435+\frac{1-0.2779}{0.4755}\approx 2.462~\text{s}.
\]

\textbf{Paso 6: calcular $T$.}

En esta construcci\'on geom\'etrica,
\[
T=t_f-L.
\]
Por tanto,
\[
T\approx 2.462-0.359=2.103~\text{s}.
\]

\textbf{Conclusi\'on del ejemplo.}

Este procedimiento es muy importante porque muchos m\'etodos de sintonizaci\'on industrial no parten del modelo exacto, sino de una curva experimental. Saber pasar de una curva a un modelo FOPDT es una habilidad central en control de procesos.

\section{Controladores PID}
\subsection{Por qu\'e el PID es tan importante}
El PID es el controlador m\'as usado en la industria porque combina:
\begin{itemize}
\item rapidez de respuesta,
\item eliminaci\'on de error estacionario,
\item amortiguamiento del transitorio.
\end{itemize}

\subsection{Forma ideal}
La forma ideal del PID es
\begin{equation}
G_c(s)=K_c\left(1+\frac{1}{\tau_i s}+\tau_d s\right).
\label{eq:pidideal}
\end{equation}

En tiempo:
\begin{equation}
u(t)=K_c e(t)+\frac{K_c}{\tau_i}\int_0^t e(\xi)\,d\xi+K_c\tau_d \frac{de(t)}{dt}.
\end{equation}

\subsection{Forma paralela}
Otra forma muy com\'un es
\begin{equation}
G_c(s)=K_p+\frac{K_i}{s}+K_d s,
\end{equation}
con las relaciones
\begin{equation}
K_p=K_c,\qquad K_i=\frac{K_c}{\tau_i},\qquad K_d=K_c\tau_d.
\end{equation}

\subsection{Ejemplo guiado: pasar de PID ideal a forma paralela}
Supongamos que ya se encontr\'o un PID ideal con
\[
K_c=7.56,\qquad \tau_i=0.76195,\qquad \tau_d=0.19049.
\]

\textbf{Paso 1: escribir la forma ideal.}

\[
G_c(s)=7.56\left(1+\frac{1}{0.76195s}+0.19049s\right).
\]

\textbf{Paso 2: usar las relaciones entre par\'ametros.}

Sabemos que
\[
K_p=K_c,\qquad K_i=\frac{K_c}{\tau_i},\qquad K_d=K_c\tau_d.
\]

\textbf{Paso 3: calcular cada constante.}

\[
K_p=7.56.
\]

\[
K_i=\frac{7.56}{0.76195}\approx 9.922.
\]

\[
K_d=7.56(0.19049)\approx 1.440.
\]

\textbf{Paso 4: escribir la forma paralela.}

\[
G_c(s)=7.56+\frac{9.922}{s}+1.440\,s.
\]

\textbf{Por qu\'e este ejemplo es importante.}

Porque en la pr\'actica muchas tablas dan el controlador en forma ideal, pero MATLAB, Simulink o ciertos bloques industriales piden $K_p$, $K_i$ y $K_d$. Si no sabes convertir entre formas, puedes implementar mal un controlador aunque la sintonizaci\'on te\'orica sea correcta.

\subsection{Forma serie}
Tambi\'en se usa
\begin{equation}
G_c(s)=K_c^S\left(1+\frac{1}{\tau_i^S s}\right)\left(1+\tau_d^S s\right).
\end{equation}

La conversi\'on de serie a ideal es
\begin{align}
K_c^I &= K_c^S\left(1+\frac{\tau_d^S}{\tau_i^S}\right),\\
\tau_i^I &= \tau_i^S+\tau_d^S,\\
\tau_d^I &= \left(\frac{1}{\tau_d^S}+\frac{1}{\tau_i^S}\right)^{-1}.
\end{align}

El paso de ideal a serie no siempre es posible. Una condici\'on cl\'asica es
\begin{equation}
\tau_i^I \ge 4\tau_d^I.
\end{equation}

\subsection{Efecto de cada t\'ermino}
\subsubsection{Acci\'on proporcional}
\[
u_P(t)=K_c e(t).
\]
Efectos t\'ipicos:
\begin{itemize}
\item disminuye el tiempo de subida,
\item reduce el error estacionario, pero no siempre lo elimina,
\item puede aumentar el sobreimpulso,
\item puede degradar la estabilidad si se exagera.
\end{itemize}

\subsubsection{Acci\'on integral}
\[
u_I(t)=\frac{K_c}{\tau_i}\int_0^t e(\xi)d\xi.
\]
Efectos:
\begin{itemize}
\item elimina el error estacionario ante escal\'on,
\item tiende a hacer la respuesta m\'as oscilatoria,
\item puede aumentar sobreimpulso y tiempo de establecimiento.
\end{itemize}

\subsubsection{Acci\'on derivativa}
\[
u_D(t)=K_c\tau_d \dot{e}(t).
\]
Efectos:
\begin{itemize}
\item anticipa cambios del error,
\item disminuye sobreimpulso,
\item mejora amortiguamiento,
\item es sensible al ruido.
\end{itemize}

\subsection{Resumen cualitativo}
Las diapositivas resumen bien esto:
\begin{itemize}
\item aumentar $K_c$: baja $t_r$, sube $M_p$, puede empeorar estabilidad;
\item aumentar $\tau_d$ o el peso derivativo: baja $M_p$ y mejora amortiguamiento;
\item aumentar la acci\'on integral: reduce mucho error estacionario, pero vuelve m\'as delicado el sistema.
\end{itemize}

\subsection{Proportional kick}
En un PI o PID paralelo, un cambio abrupto de referencia puede producir un salto fuerte en la salida del controlador porque la parte proporcional reacciona instant\'aneamente al error.

\subsection{Integrator windup}
Si el actuador satura, el integrador puede seguir acumulando error y luego tardar mucho en recuperarse. Esto produce respuestas lentas, grandes sobreimpulsos y desempe\~no pobre.

La soluci\'on t\'ipica es usar \textbf{anti-windup}.

\section{Clasificaci\'on de m\'etodos de sintonizaci\'on}
Seg\'un las clases, los m\'etodos de sintonizaci\'on pueden clasificarse as\'i:
\begin{enumerate}[label=\arabic*.]
\item Basados en respuesta escal\'on.
\item Basados en minimizar un criterio de desempe\~no.
\item Basados en respuestas deseadas de lazo cerrado.
\item Basados en robustez y estabilidad robusta.
\item Basados en frecuencia \'ultima o \emph{ultimate cycling}.
\item Otros m\'etodos pr\'acticos o semiemp\'iricos.
\end{enumerate}

\section{M\'etodo de Ziegler-Nichols}
\subsection{Versi\'on de malla abierta}
Se parte de un modelo aproximado
\[
G_p(s)\approx \frac{Ke^{-Ls}}{Ts+1}.
\]

Para PID ideal, una regla cl\'asica es:
\begin{align}
K_c &= 1.2\frac{T}{KL},\\
\tau_i &= 2L,\\
\tau_d &= 0.5L.
\end{align}

\subsection{Versi\'on de malla cerrada}
Se aumenta la ganancia proporcional hasta que el sistema quede en oscilaci\'on sostenida. Eso da:
\begin{itemize}
\item $K_u$: ganancia \'ultima o cr\'itica,
\item $P_u$: periodo de oscilaci\'on correspondiente.
\end{itemize}

Para PID ideal:
\begin{align}
K_c &= 0.6K_u,\\
\tau_i &= \frac{P_u}{2},\\
\tau_d &= \frac{P_u}{8}.
\end{align}

\subsection{Observaci\'on importante}
Dependiendo de la forma de PID usada por el profesor o el software (ideal, serie, paralelo, interactuante), los n\'umeros pueden cambiar un poco. Por eso siempre hay que verificar qu\'e estructura est\'a usando la tabla.

\section{M\'etodo de Cohen-Coon}
Tambi\'en usa un modelo FOPDT
\[
G_p(s)\approx \frac{Ke^{-Ls}}{Ts+1},
\]
pero busca una respuesta menos limitada que Ziegler-Nichols en algunos casos.

Para PID, las f\'ormulas de las diapositivas son
\begin{align}
K_c &= \frac{1}{K}\frac{T}{L}\left(\frac{16T+3L}{12T}\right),\\
\tau_i &= L\left(\frac{32+6(L/T)}{13+8(L/T)}\right),\\
\tau_d &= \frac{4L}{11+2(L/T)}.
\end{align}

Ventaja:
\begin{itemize}
\item puede ser m\'as agresivo y flexible.
\end{itemize}

Desventaja:
\begin{itemize}
\item tiende a ser menos robusto; peque\~nos cambios del proceso pueden empeorar mucho el cierre de lazo.
\end{itemize}

\section{S\'intesis directa e IMC}
\subsection{Idea}
En vez de elegir directamente $K_c$, $\tau_i$ y $\tau_d$, se elige primero c\'omo se quiere que sea la funci\'on de transferencia en lazo cerrado $T(s)$.

Si el sistema es de realimentaci\'on unitaria:
\begin{equation}
T(s)=\frac{G_c(s)G_p(s)}{1+G_c(s)G_p(s)}.
\end{equation}

Despejando el controlador:
\begin{equation}
G_c(s)=\frac{T(s)}{G_p(s)\left(1-T(s)\right)}.
\label{eq:sintesisdirecta}
\end{equation}

Esa es la ecuaci\'on de s\'intesis.

\subsection{Advertencia}
No siempre se puede cancelar libremente cualquier din\'amica del proceso. Si el proceso tiene ceros de fase no m\'inima o retardos, cancelarlos puede generar un controlador inestable o no realizable.

\section{Ejemplo resuelto: s\'intesis directa con fase no m\'inima}
En las diapositivas aparece el proceso
\begin{equation}
G_p(s)=\frac{1-9s}{(15s+1)(3s+1)}.
\end{equation}

\textbf{Paso 1: proponer una funci\'on de transferencia deseada para el lazo cerrado.}

Si se intenta imponer
\begin{equation}
T(s)=\frac{1}{\lambda s+1},
\end{equation}
entonces por \eqref{eq:sintesisdirecta}:
\[
G_c(s)=\frac{\frac{1}{\lambda s+1}}{\frac{1-9s}{(15s+1)(3s+1)}\left(1-\frac{1}{\lambda s+1}\right)}.
\]

Simplificando:
\begin{equation}
G_c(s)= -\frac{(15s+1)(3s+1)}{\lambda s(9s-1)}.
\end{equation}

\textbf{Paso 2: revisar si el controlador tiene sentido f\'isico.}

El problema es que aparece un polo en
\[
s=\frac{1}{9},
\]
que est\'a en el semiplano derecho. Es decir, el controlador ser\'ia inestable.

\textbf{Por qu\'e pasa esto.}

Porque el proceso tiene un cero de fase no m\'inima en $s=1/9$. Al intentar imponer un lazo cerrado demasiado simple, el dise\~no termina tratando de cancelar una din\'amica que no conviene cancelar.

\subsection{Correcci\'on del objetivo de lazo cerrado}
Para evitar cancelar el cero de fase no m\'inima, se propone
\begin{equation}
T(s)=\frac{1-9s}{(\lambda s+1)^2}.
\end{equation}

Entonces:
\[
G_c(s)=\frac{T(s)}{G_p(s)(1-T(s))}
\]
lleva a
\begin{equation}
G_c(s)=\frac{45s^2+18s+1}{s(\lambda^2 s+2\lambda+9)}.
\end{equation}

\textbf{Paso 3: interpretar la nueva forma.}

Ahora el cero no m\'inimo del proceso se conserva en el lazo cerrado deseado. Eso evita introducir un polo inestable en el controlador.

\subsection{Escritura como PID ideal en serie con filtro}
Queremos escribirlo como
\begin{equation}
G_c(s)=K_c\left(\frac{\tau_i\tau_d s^2+\tau_i s+1}{\tau_i s}\right)\left(\frac{1}{\tau_f s+1}\right).
\end{equation}

Comparando coeficientes:
\begin{align}
\tau_i &= 18,\\
\tau_d &= 2.5,\\
K_c &= \frac{18}{2\lambda+9},\\
\tau_f &= \frac{\lambda^2}{2\lambda+9}.
\end{align}

\textbf{Paso 4: entender por qu\'e aparece un filtro.}

El filtro aparece porque el controlador no es solo un PID ideal puro; adem\'as necesita una din\'amica adicional para que la realizaci\'on sea consistente con el objetivo de lazo cerrado.

Este ejemplo es muy importante porque ense\~na una idea profunda del curso:
\begin{itemize}
\item no se debe cancelar un cero de fase no m\'inima,
\item el objetivo de lazo cerrado debe reformularse para respetar la estructura del proceso.
\end{itemize}

\section{Ejemplo resuelto: sintonizaci\'on por Ziegler-Nichols de malla cerrada}
Retomemos el proceso
\[
G_p(s)=\frac{1}{(s+1)(0.5s+1)(0.2s+1)}.
\]

Ya se obtuvo:
\[
K_u=12.6,\qquad P_u=1.5239~\text{s}.
\]

\textbf{Paso 1: identificar qu\'e datos requiere este m\'etodo.}

La versi\'on cerrada de Ziegler-Nichols no usa $K$, $L$ y $T$. Usa la ganancia \'ultima $K_u$ y el per\'iodo cr\'itico $P_u$, que se obtienen llevando el lazo al borde de la estabilidad.

\textbf{Paso 2: escribir la regla del m\'etodo.}

Para un PID ideal:
\[
K_c=0.6K_u,\qquad \tau_i=\frac{P_u}{2},\qquad \tau_d=\frac{P_u}{8}.
\]

\textbf{Paso 3: sustituir valores.}

Aplicando Ziegler-Nichols de malla cerrada para PID ideal:
\begin{align}
K_c &= 0.6K_u = 0.6(12.6)=7.56,\\
\tau_i &= \frac{P_u}{2}=0.76195~\text{s},\\
\tau_d &= \frac{P_u}{8}=0.19049~\text{s}.
\end{align}

\textbf{Paso 4: escribir el controlador resultante.}

Por tanto, el PID ideal queda
\begin{equation}
G_c(s)=7.56\left(1+\frac{1}{0.76195s}+0.19049s\right).
\end{equation}

\textbf{Paso 5: convertir a forma paralela si se necesita.}

En forma paralela:
\begin{align}
K_p &= 7.56,\\
K_i &= \frac{7.56}{0.76195}\approx 9.922,\\
K_d &= 7.56(0.19049)\approx 1.440.
\end{align}

\textbf{Conclusi\'on del ejemplo.}

Este m\'etodo es muy r\'apido, pero tambi\'en suele ser agresivo. Por eso se usa como punto de partida y no como punto final incuestionable.

\section{Ejemplo resuelto: identificaci\'on por 63\% y sinton\'ia abierta}
Usando la aproximaci\'on del proceso de tercer orden por FOPDT:
\[
K\approx 1,\qquad L\approx 0.359,\qquad T\approx 2.103.
\]

\subsection{Ziegler-Nichols de malla abierta}
\textbf{Paso 1: recordar qu\'e datos necesita el m\'etodo.}

La versi\'on abierta de Ziegler-Nichols se basa en el modelo
\[
G_p(s)\approx \frac{Ke^{-Ls}}{Ts+1}.
\]
Por eso se parte de $K$, $L$ y $T$ estimados desde la respuesta al escal\'on.

\textbf{Paso 2: escribir la regla para PID.}

\[
K_c=1.2\frac{T}{KL},\qquad \tau_i=2L,\qquad \tau_d=0.5L.
\]

\textbf{Paso 3: sustituir valores.}

Para PID:
\begin{align}
K_c &= 1.2\frac{T}{KL}=1.2\frac{2.103}{1(0.359)}\approx 7.03,\\
\tau_i &= 2L\approx 0.718,\\
\tau_d &= 0.5L\approx 0.180.
\end{align}

\textbf{Interpretaci\'on.}

Este resultado depende por completo de qu\'e tan buena fue la identificaci\'on de $L$ y $T$. Si la tangente se traza mal, la sinton\'ia tambi\'en sale mal.

\subsection{Cohen-Coon}
\textbf{Paso 1: formar la relaci\'on $r=L/T$.}

Primero se calcula
\[
r=\frac{L}{T}\approx \frac{0.359}{2.103}\approx 0.171.
\]

\textbf{Por qu\'e aparece esta raz\'on.}

Porque Cohen-Coon corrige la agresividad del controlador seg\'un la importancia relativa del retardo respecto a la din\'amica lenta del proceso.

\textbf{Paso 2: sustituir en las ecuaciones del m\'etodo.}

Entonces:
\begin{align}
K_c &=
\frac{1}{K}\frac{T}{L}\left(\frac{16T+3L}{12T}\right)
\approx 8.06,\\
\tau_i &=L\left(\frac{32+6r}{13+8r}\right)\approx 0.825,\\
\tau_d &=\frac{4L}{11+2r}\approx 0.127.
\end{align}

\subsection{Comparaci\'on}
\begin{itemize}
\item ZN suele producir respuestas m\'as agresivas.
\item Cohen-Coon tambi\'en puede ser agresivo, y a veces menos robusto.
\item Ninguno de los dos sustituye la validaci\'on posterior por simulaci\'on.
\end{itemize}

\textbf{Conclusi\'on del ejemplo.}

Lo importante aqu\'i no es memorizar n\'umeros, sino entender el flujo completo:
\begin{enumerate}[label=\arabic*.]
\item medir una respuesta temporal;
\item aproximarla por un modelo simple;
\item aplicar una regla de sintonizaci\'on;
\item verificar por simulaci\'on si el resultado realmente funciona.
\end{enumerate}

\section{Interpretaci\'on de simulaciones PID}
Cuando en una simulaci\'on se comparan $y(t)$ y $u(t)$, hay varias cosas que conviene mirar:
\begin{itemize}
\item si $y(t)$ llega al valor deseado,
\item si hay o no sobreimpulso,
\item cu\'anto tarda en establecerse,
\item si $u(t)$ es muy grande o muy brusca,
\item si hay saturaci\'on del actuador,
\item si el derivativo introduce picos fuertes.
\end{itemize}

Es normal que la se\~nal de control obtenida por una reconstrucci\'on manual en MATLAB no coincida exactamente con la se\~nal de un bloque PID de Simulink:
\begin{itemize}
\item Simulink suele incluir filtro derivativo,
\item puede haber diferencias de solver y discretizaci\'on,
\item la implementaci\'on del derivativo puede actuar sobre el error o sobre la medici\'on.
\end{itemize}

\section{Procedimiento general para resolver ejercicios de examen}
\subsection{Si te piden modelar}
\begin{enumerate}[label=\arabic*.]
\item Escribe la ecuaci\'on diferencial.
\item Aplica Laplace con condiciones iniciales nulas.
\item Despeja $Y(s)/U(s)$.
\item Identifica polos, ceros y ganancia.
\end{enumerate}

\subsection{Si te piden estabilidad}
\begin{enumerate}[label=\arabic*.]
\item Escribe el polinomio caracter\'istico.
\item Verifica si todos los coeficientes tienen signo adecuado.
\item Usa Routh-Hurwitz.
\item Si hay par\'ametros, imp\'on positividad en la primera columna.
\end{enumerate}

\subsection{Si te piden m\'etricas temporales}
\begin{enumerate}[label=\arabic*.]
\item Lleva el sistema a forma est\'andar si es posible.
\item Identifica $\zeta$ y $\omega_n$.
\item Usa las f\'ormulas de $t_p$, $M_p$, $t_{ss}$ y $t_r$.
\end{enumerate}

\subsection{Si te piden sintonizaci\'on PID}
\begin{enumerate}[label=\arabic*.]
\item Decide si el m\'etodo es abierto, cerrado o por s\'intesis.
\item Si es ZN abierto o Cohen-Coon, identifica $K$, $L$ y $T$.
\item Si es ZN cerrado, identifica $K_u$ y $P_u$.
\item Si es s\'intesis directa, define $T(s)$ y usa
\[
G_c(s)=\frac{T(s)}{G_p(s)(1-T(s))}.
\]
\item Al final, verifica que el controlador sea realizable y estable.
\end{enumerate}

\section{Resumen r\'apido de f\'ormulas}
\subsection{Primer orden}
\[
G(s)=\frac{K}{\tau s+1},\qquad
t_{ss}(2\%)\approx 4\tau.
\]

\subsection{Segundo orden}
\[
G(s)=\frac{\omega_n^2}{s^2+2\zeta\omega_n s+\omega_n^2},
\qquad
\omega_d=\omega_n\sqrt{1-\zeta^2}.
\]
\[
t_p=\frac{\pi}{\omega_d},\qquad
M_p=e^{-\pi\zeta/\sqrt{1-\zeta^2}}100\%,\qquad
t_{ss}\approx \frac{4}{\zeta\omega_n}.
\]

\subsection{PID ideal}
\[
G_c(s)=K_c\left(1+\frac{1}{\tau_i s}+\tau_d s\right).
\]
\[
K_p=K_c,\qquad
K_i=\frac{K_c}{\tau_i},\qquad
K_d=K_c\tau_d.
\]

\subsection{ZN cerrado}
\[
K_c=0.6K_u,\qquad
\tau_i=\frac{P_u}{2},\qquad
\tau_d=\frac{P_u}{8}.
\]

\subsection{ZN abierto}
\[
K_c=1.2\frac{T}{KL},\qquad
\tau_i=2L,\qquad
\tau_d=0.5L.
\]

\subsection{Cohen-Coon PID}
\[
K_c=\frac{1}{K}\frac{T}{L}\left(\frac{16T+3L}{12T}\right),
\]
\[
\tau_i=L\left(\frac{32+6(L/T)}{13+8(L/T)}\right),
\qquad
\tau_d=\frac{4L}{11+2(L/T)}.
\]

\subsection{S\'intesis directa}
\[
G_c(s)=\frac{T(s)}{G_p(s)(1-T(s))}.
\]

\section{Conclusi\'on}
Si hubiera que resumir todo el material en pocas ideas, ser\'ian estas:
\begin{enumerate}[label=\arabic*.]
\item El modelo del sistema siempre es el punto de partida.
\item Los polos determinan estabilidad y forma de la respuesta.
\item Routh-Hurwitz permite estudiar estabilidad sin calcular las ra\'ices.
\item La respuesta al escal\'on conecta matem\'aticas con intuici\'on f\'isica.
\item El PID funciona porque cada t\'ermino hace algo distinto: rapidez, precisi\'on y amortiguamiento.
\item Sintonizar no es ``probar a ver qu\'e pasa''; es aplicar una metodolog\'ia concreta.
\item Los m\'etodos de sintonizaci\'on son aproximaciones: siempre hay que validar el resultado final.
\end{enumerate}

Si usas este documento para estudiar, la mejor estrategia es:
\begin{enumerate}[label=\arabic*.]
\item leer primero la teor\'ia,
\item luego rehacer a mano los ejemplos resueltos,
\item y finalmente intentar resolverlos sin mirar el procedimiento.
\end{enumerate}

\end{document}
